\documentclass[12pt, a4paper]{ctexart}
\usepackage{fontspec}
\usepackage{xcolor}
\usepackage{hyperref}
\usepackage{geometry}
\usepackage{enumitem}
\usepackage{parskip}
\usepackage{titlesec}
\usepackage{tcolorbox}
\tcbuselibrary{breakable}
\usepackage{setspace}
\usepackage{listings}
\usepackage{needspace}

\geometry{left=2.5cm, right=2.5cm, top=2.5cm, bottom=2.5cm}

\setmainfont{Times New Roman}
\setCJKmainfont{SimSun}

\hypersetup{
    colorlinks=true,
    linkcolor=blue!70!black,
    urlcolor=cyan,
    pdftitle={Java异常-逐题精讲解义},
    pdfauthor={专升本-fifine老师}
}

\definecolor{myblue}{RGB}{30,100,200}
\definecolor{lightblue}{RGB}{230,240,255}
\definecolor{darkblue}{RGB}{0,50,120}

\newtcolorbox{bluebox}[1][]{
    colback=lightblue,
    colframe=myblue,
    coltitle=darkblue,
    fonttitle=\bfseries,
    title={#1},
    boxrule=0.8pt,
    arc=4pt,
    left=6pt,
    right=6pt,
    top=6pt,
    bottom=6pt,
    before skip=8pt,
    after skip=8pt,
    breakable
}

\usepackage{etoolbox}
\BeforeBeginEnvironment{bluebox}{\par\vfill}%
\AfterEndEnvironment{bluebox}{\par\vfill}%

\titleformat{\section}
  {\normalfont\Large\bfseries\color{myblue}}{\thesection}{1em}{}
\titlespacing*{\section}{0pt}{3.5ex plus 1ex minus .2ex}{1.5ex plus .2ex}

\title{\textcolor{myblue}{\textbf{Java异常 - 逐题精讲解义}}}
\author{专升本-fifine老师 教学组}
\date{2026年03月02日}

\lstset{
    language=Java,
    basicstyle=\ttfamily\small,
    keywordstyle=\color{blue},
    commentstyle=\color{green!60!black},
    stringstyle=\color{orange},
    numbers=left,
    numberstyle=\tiny\color{gray},
    frame=single,
    breaklines=true,
    columns=fullflexible,
    showstringspaces=false
}

\newcommand{\code}[1]{\texttt{#1}}

\begin{document}

\maketitle
\thispagestyle{empty}
\newpage

\section{异常}

\begin{enumerate}[label=\textbf{\arabic*.}]

	\item \textbf{以下哪个异常类用于处理数组下标越界异常？}
	      \begin{enumerate}[label=\Alph*.]
		      \item ArrayIndexOutOfBoundsException
		      \item NullPointerException
		      \item ArithmeticException
		      \item ClassCastException
	      \end{enumerate}

	      \begin{bluebox}{答案与深度解析}
		      \textbf{答案：A. ArrayIndexOutOfBoundsException}

		      % ======== 阶段一：核心认知框架 ========
		      \textbf{【核心认知框架】}
		      \begin{itemize}[label={}, leftmargin=*, nosep]
			      \item \textbf{题目本质}：考查Java异常体系中数组相关异常类的区分
			      \item \textbf{关键区分点}：数组越界 vs 空指针 vs 算术错误 vs 类型转换错误
			      \item \textbf{命题人意图}：测试考生对Java异常分类体系的掌握程度
		      \needspace{10cm}
		      % ======== 阶段二：选项深度拆解 ========
		      \textbf{【选项原理深度拆解】}

		      \item \textbf{A. ArrayIndexOutOfBoundsException — 正确（数组索引越界）}
		      \begin{itemize}[label={}, leftmargin=*, nosep]
			      \item \textbf{字节码铁证}：
			            \begin{lstlisting}[language=Java]
int[] arr = new int[5];
arr[5] = 10; // 运行时抛出 ArrayIndexOutOfBoundsException
      \end{lstlisting}
			            \textbf{JVM保障机制}：数组访问时JVM检查索引，超出范围抛出异常

			      \item \textbf{RuntimeException深度解析}：
			            \begin{itemize}[label={}, leftmargin=*, nosep]
				            \item 属于RuntimeException（未检查异常）
				            \item JLS §14.17 规定数组访问自动边界检查
				            \item 常见场景：循环索引误用、动态数组操作
			            \end{itemize}

			      \item \textbf{命题陷阱破解}：
			            \begin{itemize}[label={}, leftmargin=*, nosep]
				            \item 干扰项："数组长度为5，访问arr[4]会越界" → 错！arr[4]是第5个元素，合法
				            \item \textbf{必考结论}：数组索引从0开始，最大有效索引为length-1
			            \end{itemize}
		      \end{itemize}

		      \item \textbf{B. NullPointerException — 错误（空指针异常）}
		      \begin{itemize}[label={}, leftmargin=*, nosep]
			      \item \textbf{触发条件}：对象引用为null时访问成员
			      \item \textbf{字节码铁证}：
			            \begin{lstlisting}[language=Java]
String str = null;
str.length(); // 抛出 NullPointerException
      \end{lstlisting}
			      \item \textbf{核心区别}：数组越界是索引问题，空指针是引用问题
		      \end{itemize}

		      \item \textbf{C. ArithmeticException — 错误（算术异常）}
		      \begin{itemize}[label={}, leftmargin=*, nosep]
			      \item \textbf{触发条件}：算术运算错误，如除零
			      \item \textbf{典型场景}：int x = 10 / 0; // 抛出 ArithmeticException
			      \item \textbf{核心区别}：与数组操作无关
		      \end{itemize}

		      \item \textbf{D. ClassCastException — 错误（类型转换异常）}
		      \begin{itemize}[label={}, leftmargin=*, nosep]
			      \item \textbf{触发条件}：强制类型转换时类型不兼容
			      \item \textbf{典型场景}：Object obj = "hello"; Integer num = (Integer) obj;
			      \item \textbf{核心区别}：运行时类型检查失败，与数组索引无关
		      \end{itemize}

		      % ======== 阶段三：应背诵知识点 ========
		      \textbf{【应背诵知识点（命题人思维版）】}
		      \begin{enumerate}[leftmargin=*, nosep, label={}]
			      \item \textbf{数组相关异常体系}：
			            \begin{itemize}[label={}, leftmargin=*, nosep]
				            \item ArrayIndexOutOfBoundsException：索引越界
				            \item NegativeArraySizeException：数组大小为负
				            \item ArrayStoreException：数组类型存储错误
			            \end{itemize}

			      \item \textbf{RuntimeException vs CheckedException}：
			            \begin{itemize}[label={}, leftmargin=*, nosep]
				            \item RuntimeException：编译器不检查，可不处理（题中四个选项全是RuntimeException）
				            \item CheckedException：必须捕获或声明throws
			            \end{itemize}

			      \item \textbf{选项辨析速查表}：
			      \item 
			            \begin{tabular}{|c|c|c|c|}
				            \hline
				            \textbf{选项}                    & \textbf{触发场景} & \textbf{错误根源} & \textbf{记忆口诀}           \\
				            \hline
				            ArrayIndexOutOfBoundsException & 索引超范围         & 数组访问越界        & \textcolor{green}{"越界"} \\
				            \hline
				            NullPointerException           & 引用为null       & 空对象访问成员       & \textcolor{red}{"空指针"}  \\
				            \hline
				            ArithmeticException            & 除零等           & 算术运算错误        & \textcolor{red}{"算术错"}  \\
				            \hline
				            ClassCastException             & 强制转换          & 类型不兼容         & \textcolor{red}{"转换错"}  \\
				            \hline
			            \end{tabular}
		      \end{enumerate}

		      % ======== 阶段四：命题趋势与实战策略 ========
		      \textbf{【命题趋势与实战策略】}

		      \textbf{考场秒杀三步法}：
		      \begin{enumerate}[label={}, leftmargin=*, nosep]
			      \item 看到"数组下标/索引" → 锁定ArrayIndexOutOfBoundsException
			      \item 排除明显无关选项（空指针、算术错、类型转换）
			      \item 确认索引值是否超出0到length-1范围
		      \end{enumerate}
		      \end{itemize}
	      \end{bluebox}

	      \vspace{1em}

	\item \textbf{在Java中，以下哪个关键字用于捕获异常？}
	      \begin{enumerate}[label=\Alph*.]
		      \item try
		      \item catch
		      \item throw
		      \item throws
	      \end{enumerate}

	      \begin{bluebox}{答案与深度解析}
		      \textbf{答案：B. catch}

		      \textbf{【核心认知框架】}
		      \begin{itemize}[label={}, leftmargin=*, nosep]
			      \item \textbf{题目本质}：考查异常处理关键字的功能区分
			      \item \textbf{关键区分点}：try（监测）、catch（捕获处理）、throw（抛出）、throws（声明）
			      \item \textbf{命题人意图}：区分四个易混关键字的不同作用

		      \textbf{【选项原理深度拆解】}

		      \item \textbf{A. try — 错误（监测异常而非捕获）}
		      \begin{itemize}[label={}, leftmargin=*, nosep]
			      \item try块用于监测可能抛出异常的代码
			      \item 必须配合catch或finally使用
			      \item 单独try无法处理异常
		      \end{itemize}

		      \item \textbf{B. catch — 正确（捕获并处理异常）}
		      \begin{itemize}[label={}, leftmargin=*, nosep]
			      \item 紧跟try块后，用于捕获指定类型的异常
			      \item 可以有多个catch块（异常子类在前，父类在后）
			      \item 是异常处理的实际执行者
		      \end{itemize}

		      \item \textbf{C. throw — 错误（手动抛出异常）}
		      \begin{itemize}[label={}, leftmargin=*, nosep]
			      \item 用于在代码中主动抛出异常
			      \item 格式：throw new Exception("信息");
			      \item 作用在方法内部
		      \end{itemize}

		      \item \textbf{D. throws — 错误（声明异常）}
		      \begin{itemize}[label={}, leftmargin=*, nosep]
			      \item 用于方法签名，声明该方法可能抛出的异常
			      \item 格式：public void method() throws Exception \{\}
			      \item 将异常处理责任交给调用者
		      \end{itemize}

		      \textbf{【应背诵知识点（命题人思维版）】}

		      \textbf{关键字功能速查表}：
		      \begin{tabular}{|c|c|c|}
			      \hline
			      \textbf{关键字} & \textbf{作用位置} & \textbf{功能描述} \\
			      \hline
			      try          & 代码块           & 监测异常发生        \\
			      \hline
			      catch        & try后          & 捕获并处理异常       \\
			      \hline
			      throw        & 方法内           & 手动抛出异常        \\
			      \hline
			      throws       & 方法签名          & 声明方法可能抛出的异常   \\
			      \hline
		      \end{tabular}

		      \textbf{记忆口诀}：try监测catch抓，throw抛出throws声明

		      \textbf{【命题趋势与实战策略】}

		      \textbf{考场秒杀}：catch是唯一真正"捕获"异常的关键字，题干问"捕获"选catch
		      \end{itemize}

	      \end{bluebox}

	      \vspace{1em}

	\item \textbf{以下关于Java中异常处理中finally块的说法正确的是}
	      \begin{enumerate}[label=\Alph*.]
		      \item 无论是否发生异常都会执行
		      \item 只有发生异常才会执行
		      \item 只有不发生异常才会执行
		      \item 只有在try块中有return语句时才会执行
	      \end{enumerate}

	      \begin{bluebox}{答案与深度解析}
		      \textbf{答案：A. 无论是否发生异常都会执行}

		      \textbf{【核心认知框架】}
		      \begin{itemize}[label={}, leftmargin=*, nosep]
			      \item \textbf{题目本质}：考查finally块的核心特性——必然执行
			      \item \textbf{关键区分点}：finally vs try/catch的执行顺序
			      \item \textbf{命题人意图}：测试对异常处理流程的深入理解

		      \textbf{【选项原理深度拆解】}

		      \item \textbf{A. 无论是否发生异常都会执行 — 正确（finally核心特性）}
		      \begin{itemize}[label={}, leftmargin=*, nosep]
			      \item \textbf{JVM规范依据}：JVMS §2.13.5 规定finally块必须执行
			      \item \textbf{执行时机}：try/catch执行后，无论是否有异常都会执行
			      \item \textbf{即使有return}：finally仍会在return之前执行
		      \end{itemize}

		      \item \textbf{B. 只有发生异常才会执行 — 错误}
		      \begin{itemize}[label={}, leftmargin=*, nosep]
			      \item 这是对finally的常见误解
			      \item 正常情况下（无异常）try执行完也会执行finally
		      \end{itemize}

		      \item \textbf{C. 只有不发生异常才会执行 — 错误}
		      \begin{itemize}[label={}, leftmargin=*, nosep]
			      \item 与B选项相反，同样错误
			      \item finally与是否异常无关
		      \end{itemize}

		      \item \textbf{D. 只有在try块中有return语句时才会执行 — 错误}
		      \begin{itemize}[label={}, leftmargin=*, nosep]
			      \item 恰恰相反：有return时finally仍会执行
			      \item 注意：return会等待finally执行完再返回
		      \end{itemize}

		      \textbf{【应背诵知识点（命题人思维版）】}

		      \textbf{finally块执行规则}：
		      \begin{enumerate}[label={}, leftmargin=*, nosep]
			      \item 无论try是否抛出异常，finally必定执行
			      \item 即使try/catch中有return，finally仍会在return前执行
			      \item 如果JVM退出（System.exit()），finally不执行
			      \item 如果线程被杀死，finally也不执行
		      \end{enumerate}

		      \textbf{命题陷阱破解}：
		      \begin{itemize}[label={}, leftmargin=*, nosep]
			      \item 陷阱1："catch执行了finally才执行" → 错！正常流程try执行完就执行finally
			      \item 陷阱2："return会跳过finally" → 错！return会等待finally完成
		      \end{itemize}

		      \textbf{【命题趋势与实战策略】}

		      \textbf{考场秒杀}：牢记finally三字真言——"必然执行"
		      \end{itemize}

	      \end{bluebox}

	      \vspace{1em}

	\item \textbf{以下哪种情况会导致OutOfMemoryError？}
	      \begin{enumerate}[label=\Alph*.]
		      \item 数组越界
		      \item 除数为0
		      \item 内存泄漏
		      \item 空指针异常
	      \end{enumerate}

	      \begin{bluebox}{答案与深度解析}
		      \textbf{答案：C. 内存泄漏}

		      \textbf{【核心认知框架】}
		      \begin{itemize}[label={}, leftmargin=*, nosep]
			      \item \textbf{题目本质}：区分Error类型与Exception类型
			      \item \textbf{关键区分点}：OutOfMemoryError是Error而非Exception
			      \item \textbf{命题人意图}：测试对Java异常体系层次结构的理解

		      \textbf{【选项原理深度拆解】}

		      \item \textbf{A. 数组越界 — 错误（导致RuntimeException）}
		      \begin{itemize}[label={}, leftmargin=*, nosep]
			      \item 数组越界抛出 ArrayIndexOutOfBoundsException
			      \item 属于RuntimeException，不是Error
		      \end{itemize}

		      \item \textbf{B. 除数为0 — 错误（导致ArithmeticException）}
		      \begin{itemize}[label={}, leftmargin=*, nosep]
			      \item 除数为0在整数运算时抛出 ArithmeticException
			      \item 浮点数除以0得到Infinity，不会抛异常
		      \end{itemize}

		      \item \textbf{C. 内存泄漏 — 正确（导致OutOfMemoryError）}
		      \begin{itemize}[label={}, leftmargin=*, nosep]
			      \item \textbf{Error vs Exception}：Error是严重错误，通常由JVM抛出
			      \item \textbf{JVM保障机制}：当堆内存耗尽时抛出OutOfMemoryError
			      \item \textbf{常见原因}：长生命周期对象持有短生命周期引用、静态集合未清理、大对象创建
		      \end{itemize}

		      \item \textbf{D. 空指针异常 — 错误（导致NullPointerException）}
		      \begin{itemize}[label={}, leftmargin=*, nosep]
			      \item NullPointerException属于RuntimeException
			      \item 与内存无关
		      \end{itemize}

		      \textbf{【应背诵知识点（命题人思维版）】}

		      \textbf{异常体系结构}：
		      \begin{itemize}[label={}, leftmargin=*, nosep]
			      \item Throwable
			            \begin{itemize}[label={}, leftmargin=*, nosep]
				            \item Error（严重错误，程序无法处理）
				                  \begin{itemize}[label={}, leftmargin=*, nosep]
					                  \item OutOfMemoryError
					                  \item StackOverflowError
				                  \end{itemize}
				            \item Exception（可处理异常）
				                  \begin{itemize}[label={}, leftmargin=*, nosep]
					                  \item RuntimeException（未检查异常）
					                  \item CheckedException（已检查异常）
				                  \end{itemize}
			            \end{itemize}
		      \end{itemize}

		      \textbf{OutOfMemoryError常见类型}：
		      \begin{itemize}[label={}, leftmargin=*, nosep]
			      \item Java heap space：堆内存不足
			      \item GC overhead limit exceeded：GC开销过大
			      \item Unable to create new native thread：线程数超限
			      \item Metaspace：元空间不足（Java 8+）
		      \end{itemize}

		      \textbf{【命题趋势与实战策略】}

		      \textbf{考场秒杀}：OutOfMemoryError=内存问题，选"内存泄漏"
		      \end{itemize}

	      \end{bluebox}

\end{enumerate}

\end{document}
